\section*{Chapter \ref{ch:iv1:regression} --- Regression}
\markboth{\protect\omTocMark{}{Answers: Regression}}{\protect\omTocMark{}{Answers: Regression}}

\begin{solution}{iq:iv1:regression:1}
The reverse slope can be anything above 0 and up to $1/2$. A reverse slope of $0.18$ means a correlation of $0.6$.

With $b_1 = \mathrm{Cov}(x,y)/\mathrm{Var}(x)$ and $b_2 = \mathrm{Cov}(x,y)/\mathrm{Var}(y)$, the product is
\[
b_1 b_2 = \frac{\mathrm{Cov}(x,y)^2}{\mathrm{Var}(x)\mathrm{Var}(y)} = \rho^2 .
\]
So $b_2 = \rho^2/2$, with $\rho^2$ in $(0, 1]$. The two slopes always have the same sign. If $b_2 = 0.18$ then $\rho^2 = 0.36$, and the correlation is $0.6$, positive because the slopes are.

\iqwrong{``$b_2 = 1/b_1 = 1/2$.'' That holds only with $|\rho| = 1$. Each regression minimises vertical errors along its own axis, so the two lines differ, and they coincide only when all the points lie on a line.}
\iqwhy{The choice of which variable to regress on which is not free when you set a hedge ratio. The two give different ratios, and the correct one depends on what you want to minimise.}
\end{solution}

\begin{solution}{iq:iv1:regression:2}
The coefficients and $R^2$ do not change. The standard errors shrink by a factor of about $\sqrt2$, and the $t$-statistics rise by about $\sqrt2$.

$X^\top X$ and $X^\top y$ both double, so $\hat\beta = (X^\top X)^{-1}X^\top y$ is unchanged. The residuals repeat, so RSS doubles while the degrees of freedom almost double, and $s^2$ is about the same. Since $(X^\top X)^{-1}$ halves, the variance estimate halves. Exactly, with $n$ rows and $p$ coefficients, the standard errors scale by $\sqrt{(n-p)/(2n-p)}$, which for $n = 250$ and $p = 2$ is $0.706$, next to $1/\sqrt2 = 0.707$.

\iqwhy{Nothing has been learned, yet the $t$-statistic has grown. Overlapping data do the same thing: 20-day returns sampled daily look like 20 times the data, while holding about one independent observation per 20 days. With a slow signal the $t$-statistic is overstated by about $\sqrt{20} \approx 4.5$. Use non-overlapping samples or standard errors that allow for the overlap.}
\end{solution}

\begin{solution}{iq:iv1:regression:3}
Short $1.6$ dollars of the index per dollar of stock. The hedged position has a volatility of 18\%. Using the inverted reverse regression, you would short $2.5$ dollars and leave $22.5\%$.

The variance of stock minus $h$ times index is a quadratic in $h$, minimised at $h = \mathrm{Cov}/\mathrm{Var}(\text{index}) = \rho\,\sigma_s/\sigma_i = 0.8 \cdot 30/15 = 1.6$. This is the slope of the stock regressed on the index. The residual volatility is $\sigma_s\sqrt{1 - \rho^2} = 30 \cdot 0.6 = 18\%$, so the hedge removes $\rho^2 = 64\%$ of the variance.

The regression of the index on the stock has slope $\rho\,\sigma_i/\sigma_s = 0.4$, and inverting it gives $2.5$. The variance is then $30^2 + 2.5^2 \cdot 15^2 - 2 \cdot 2.5 \cdot 360 = 506.25$, a volatility of $22.5\%$. It over-hedges by the factor $1/\rho^2$.

\iqwhy{This is a daily decision on a desk. Note also that an estimated ratio on noisy returns tends to come out low, not high, because noise in the regressor shrinks the slope towards zero.}
\end{solution}

\begin{solution}{iq:iv1:regression:4}
With an intercept, $R^2$ lies in $[0, 1]$, and $R^2 = 1$ means that every residual is zero. Adding a regressor can never lower it. Without an intercept, that formula can be far below zero.

With an intercept, the total sum of squares splits into explained plus residual, because the residuals are orthogonal to the fitted values and sum to zero\index{normal equations}. So $R^2$ is a ratio between 0 and 1, and it equals the squared correlation of $y$ with the fitted values. A new regressor can reproduce the old fit by taking a zero coefficient, so RSS cannot rise.

Without an intercept, the fit is not obliged to do as well as the mean. If $y$ has a large mean and $x$ is centred on zero, RSS is close to $\sum y_i^2$, which is much larger than $\sum(y_i - \bar y)^2$, and the formula gives a large negative number.

\iqwrong{Reading a higher $R^2$ from an added regressor as evidence. It rises even for pure noise.}
\iqfollow{And out of sample?}{It can be negative for any model, when the model predicts worse than the mean.}
\end{solution}

\begin{solution}{iq:iv1:regression:5}
The $t$-statistic is $0.4$, and you would need 100 years to reach a $t$-statistic of 2.

The standard error of the intercept is about the residual volatility divided by the square root of the number of years, $10\%/\sqrt4 = 5\%$ a year. So $t = 2/5 = 0.4$. In terms of the information ratio\index{information ratio}, the alpha divided by residual volatility, here $0.2$, you have $t = \text{ratio} \cdot \sqrt{\text{years}}$. To reach 2 requires $\sqrt{\text{years}} = 2/0.2 = 10$.

\iqfollow{What if the information ratio were 1?}{Then $t = 2$ takes 4 years.}
\iqwhy{An alpha of this size cannot be confirmed from track record in a working lifetime. In practice you lean on the economic logic of the strategy, on a much higher frequency of independent bets, or on cross-sectional evidence. Performance alone is not enough.}
\end{solution}

\enlargethispage{-6\baselineskip}%
\begin{solution}{iq:iv1:regression:6}
The $t$-statistic is about $3.54$. The signal may be valuable, with an annual Sharpe ratio of about $1.1$.

A simple regression has the exact identity
\[
t^2 = \frac{R^2\,(n - 2)}{1 - R^2},
\]
so $t = \sqrt{0.005 \cdot 2498/0.995} \approx 3.54$. A low $R^2$ with a high $t$ is normal: $t$ grows with the sample, and $R^2$ does not.

The correlation with the return is $\sqrt{0.005} \approx 0.071$. Holding a position proportional to the standardised signal $x$ earns $xy$, with mean about $\rho\,\sigma_y$ and standard deviation about $\sigma_y$, so the daily Sharpe ratio is about $\rho$. Annualised, $0.071 \cdot \sqrt{252} \approx 1.12$. The same answer comes from $t/\sqrt{\text{years}} = 3.54/\sqrt{9.9}$.

\iqwrong{``An $R^2$ of half a percent means it explains nothing.'' On daily returns that is a strong signal, before costs.}
\iqwhy{Judge return prediction by Sharpe ratio after costs. This calculation also assumes the data did not come out of a search over many signals, and that is usually what you should check next.}
\end{solution}

\begin{solution}{iq:iv1:regression:7}
The estimator is $\hat\beta = (X^\top X)^{-1}X^\top y$, and it exists uniquely when $X$ has full column rank.

Minimise $L(\beta) = \lVert y - X\beta\rVert^2 = y^\top y - 2\beta^\top X^\top y + \beta^\top X^\top X\beta$. Setting the gradient $-2X^\top y + 2X^\top X\beta$ to zero gives the normal equations\index{normal equations} $X^\top X\beta = X^\top y$. The Hessian $2X^\top X$ is positive semidefinite, so this is a minimum, and it is unique if $X^\top X$ is invertible, which is the same as full column rank. With fewer rows than columns, or with a repeated column, there are many solutions.

Geometrically, $X\hat\beta$ is the orthogonal projection of $y$ onto the column space of $X$. So the residuals are orthogonal to every column, and with an intercept they sum to zero.

Do not invert. The condition number of $X^\top X$ is the square of that of $X$, so inverting loses twice as many digits. Solve by QR on $X$, or by Cholesky on $X^\top X$ if $X$ is well conditioned.

If $X$ does not fit in memory, read it in chunks and accumulate the $p \times p$ matrix $X^\top X$ and the vector $X^\top y$. They are sums over rows, so this gives the exact answer with memory proportional to $p^2$.

\iqfollow{And when rows arrive one at a time?}{Update the inverse with recursive least squares, and add a forgetting factor if the relation drifts.}
\end{solution}

\begin{solution}{iq:iv1:regression:8}
Procedure B gives the full-regression coefficient exactly. Procedure A gives $0.36$ here, not 1.

This is the Frisch--Waugh--Lovell theorem\index{Frisch--Waugh--Lovell theorem}: the coefficient of $x$ in the full regression equals the slope of the residuals of $y$ on the residuals of $x$, after removing $z$ from both. Removing $z$ from $y$ only leaves $x$ still carrying its $z$ part, which is unrelated to the residuals of $y$ and dilutes the slope. Writing $\tilde x$ for the part of $x$ not explained by $z$,
\[
\text{slope of A} = \beta \cdot \frac{\mathrm{Var}(\tilde x)}{\mathrm{Var}(x)} = \beta\,(1 - \rho^2) = 1 \cdot (1 - 0.64) = 0.36 .
\]
\iqwrong{``Partial out the market from returns, and then test the signal on the residuals.'' A common workflow, and it is procedure A. It biases the signal's coefficient towards zero by $1 - \rho^2$, and worse if the signal is highly correlated with the factor.}
\iqwhy{Neutralise the signal as well as the returns. If you sector-neutralise only the returns and your signal has a sector tilt, you understate what the signal would earn after you hedge it.}
\end{solution}

\begin{solution}{iq:iv1:regression:9}
The slope is $-0.5$. The sign is wrong: the true coefficient of $x$ is $+1$.

This is omitted-variable bias\index{omitted-variable bias}. The slope of the short regression is $\mathrm{Cov}(x, y)/\mathrm{Var}(x)$. Substituting the true model,
\[
\frac{\mathrm{Cov}(x,\, x - 3z + \varepsilon)}{\mathrm{Var}(x)} = 1 - 3 \cdot \frac{\mathrm{Cov}(x, z)}{\mathrm{Var}(x)} = 1 - 3 \cdot 0.5 = -0.5 .
\]
In general, the short slope is the true coefficient, plus the coefficient of $z$ multiplied by the slope of $z$ on $x$. It disappears only if $z$ has no effect, or if $x$ and $z$ are uncorrelated, and it does not shrink with more data.

\iqfollow{And if $x$ and $z$ are uncorrelated?}{The slope is 1, unbiased. The omitted $z$ only adds to the noise.}
\iqwhy{A signal correlated with market beta or sector will pick up their returns. Before believing a signal, regress it on the known factors and look at what remains.}
\end{solution}

\begin{solution}{iq:iv1:regression:10}
Noise in $x$ gives a slope of $1/2$. Noise in $y$ alone leaves the slope at 1. With noise in both, the true slope lies between the slope of $y$ on $x$ and the inverse of the slope of $x$ on $y$.

This is errors-in-variables\index{errors in variables}. Noise in the regressor raises $\mathrm{Var}(x)$ and leaves the covariance alone:
\[
\text{slope} = \frac{\mathrm{Cov}(x^* + n,\, y)}{\mathrm{Var}(x^*) + \mathrm{Var}(n)} = \frac{1}{1 + 1} = \frac12 .
\]
The factor $\mathrm{Var}(x^*)/\mathrm{Var}(x)$ is the attenuation. Noise in $y$ goes into the error term: it widens the standard errors but causes no bias.

With noise in both, the slope of $y$ on $x$ is too flat and the inverse of the slope of $x$ on $y$ is too steep. With the numbers above and unit noise in $y$ as well, these are $0.5$ and $2$, around the true value of 1.

\iqfollow{Is there a single line that treats $x$ and $y$ alike?}{Total least squares\index{total least squares} minimises perpendicular distance and takes the first principal direction. Its slope is
$\bigl[(s_{yy} - s_{xx}) + \sqrt{(s_{yy} - s_{xx})^2 + 4s_{xy}^2}\,\bigr]/(2s_{xy})$, here exactly 1. It assumes equal noise in the two variables and depends on their units, so standardise first or use the known noise ratio.}
\iqwhy{Estimated betas, volatilities and signals are all noisy, so regressions on them come out flat. If you size a hedge from a noisy regressor, you under-hedge.}
\end{solution}

\begin{solution}{iq:iv1:regression:11}
Each standard error is larger by a factor of about $2.29$. The sum of the coefficients is still well estimated, and yes, both $t$-statistics can be small while $F$ is large.

The variance inflation factor\index{variance inflation factor} is $1/(1 - \rho^2) = 1/0.19 = 5.26$, and its square root is $2.29$.

For unit-variance regressors, the covariance of the coefficients is proportional to the inverse of $\bigl(\begin{smallmatrix}1 & \rho\\ \rho & 1\end{smallmatrix}\bigr)$. That gives $\mathrm{Var}(b_1 + b_2) \propto 2/(1 + \rho) = 1.05$ and $\mathrm{Var}(b_1 - b_2) \propto 2/(1 - \rho) = 20$. The data pin down the sum 19 times more tightly than the difference, because the two regressors nearly always move together. The difficulty is in dividing the credit between them, which shows up as coefficients that flip sign between samples.

The $F$-test looks at the joint effect, so it does not suffer. With 100 observations and $R^2 = 0.2$, the test statistic is $F = (0.2/2)/(0.8/97) \approx 12.1$, far above the 5\% critical value of about $3.1$. In a simulation with correlation $0.95$ and true coefficients of $0.2$ each, over half the samples have both $|t| < 2$, and most of these still have a significant $F$.

\iqwhy{Drop one, combine them, or use ridge. If you only need predictions and the correlation is stable, you may do nothing, as collinearity does not harm the fit.}
\end{solution}

\begin{solution}{iq:iv1:regression:12}
The two lines are $y = x/2$ and $x = (1 + y)/2$. The correlation is $1/2$.

Because the region gives simple conditional means, you can write the regression directly. Given $x$, the height $y$ is uniform on $[0, x]$, so $E[y \mid x] = x/2$. Given $y$, the value $x$ is uniform on $[y, 1]$, so $E[x \mid y] = (1 + y)/2$. Both are straight lines, and so each is also the population regression line\index{conditional expectation}, as shown in \cref{fig:iv1:triangle}. Both slopes are $1/2$, and the product of slopes equals $\rho^2$, so $\rho = 1/2$.

\begin{omfigure}
\begin{tikzpicture}[x=4.2cm,y=4.2cm]
\fill[omDef!10] (0,0) -- (1,0) -- (1,1) -- cycle;
\draw[thick,omDef] (0,0) -- (1,0) -- (1,1) -- cycle;
\draw[->] (0,0) -- (1.2,0) node[right,font=\footnotesize] {$x$};
\draw[->] (0,0) -- (0,1.15) node[above,font=\footnotesize] {$y$};
\draw[very thick,omThm] (0,0) -- (1,0.5);
\draw[very thick,omExo] (0.5,0) -- (1,1);
\fill (2/3,1/3) circle (0.012);
\node[font=\footnotesize,anchor=north west] at (1.02,0.5) {$y = x/2$};
\node[font=\footnotesize,anchor=west] at (1.02,1) {$x = (1+y)/2$};
\node[font=\footnotesize,anchor=north west] at (0.69,0.32) {mean};
\node[below,font=\footnotesize] at (1,0) {$1$};
\node[left,font=\footnotesize] at (0,1) {$1$};
\end{tikzpicture}

\omcaption{The triangle and its two regression lines. They cross at the mean point $(2/3, 1/3)$. Drawn in the same axes, one line has slope $1/2$ and the other slope 2, on either side of the slope 1 of the diagonal edge.}\label{fig:iv1:triangle}
\end{omfigure}

As a check by integration, $E[x] = 2/3$, $E[y] = 1/3$, both variances are $1/18$ and the covariance is $1/36$. So each slope is $(1/36)/(1/18) = 1/2$.

\iqwrong{Taking the diagonal edge, $y = x$, as the line of the data. It is the boundary of the region, and not a regression line.}
\end{solution}

\begin{solution}{iq:iv1:regression:13}
OLS has no unique solution, since only the sum of the two coefficients is identified. Ridge splits the coefficient equally. Lasso does not care how it splits, so the solver decides.

For OLS, $X^\top X$ is singular, and any pair with the right sum fits equally well. A pseudo-inverse returns the minimum-norm solution, which splits equally, and a plain inverse fails or returns huge numbers of opposite sign.

For ridge, the penalty $b_1^2 + b_2^2$ for a given sum is smallest when they are equal. Each coefficient is $x^\top y/(2x^\top x + \lambda)$, and their sum is $x^\top y/(x^\top x + \lambda/2)$. So doubling a column cuts its effective penalty in half.

For lasso\index{lasso}, $|b_1| + |b_2|$ is the same for every split with the same sign, so all such splits have an equal objective. Coordinate descent from zero tends to load everything on the first column and leave the second at zero. This depends on the order and the start.

\iqwhy{With nearly identical signals the behaviour is similar, with lasso's choice jumping from day to day and creating turnover with no economic meaning. Ridge or a grouping of the signals is more stable.}
\end{solution}

\enlargethispage{-2\baselineskip}%
\begin{solution}{iq:iv1:regression:14}
Lasso gives $(2, 0, 0.2)$, and ridge gives $(1.5, -0.25, 0.6)$.

With orthonormal columns, the objective separates into one problem per coefficient. Up to constants, $\tfrac12\lVert y - Xb\rVert^2 = \tfrac12\sum_j (b_j - \hat\beta_j)^2$, where $\hat\beta_j$ is the OLS coefficient.

For ridge, minimising $\tfrac12(b - \hat\beta)^2 + \tfrac{\lambda}{2}b^2$ gives $b = \hat\beta/(1 + \lambda)$: every coefficient is divided by 2.

For lasso, minimise $\tfrac12(b - \hat\beta)^2 + \lambda|b|$. If $b > 0$ the derivative is zero at $b = \hat\beta - \lambda$, which requires $\hat\beta > \lambda$. The case $b < 0$ is symmetric. If $|\hat\beta| \le \lambda$, the kink at zero is the minimum. This is soft-thresholding\index{soft-thresholding},
\[
b = \mathrm{sign}(\hat\beta)\,\bigl(|\hat\beta| - \lambda\bigr)_+ ,
\]
plotted in \cref{fig:iv1:shrink}. So $3 \to 2$, $-0.5 \to 0$, and $1.2 \to 0.2$.

\begin{omfigure}
\begin{tikzpicture}
\begin{axis}[width=8cm,height=5.6cm,xmin=-3.3,xmax=3.3,ymin=-3.3,ymax=3.3,axis lines=middle,
  xlabel={\small OLS},ylabel={\small estimate},xlabel style={anchor=west},ylabel style={anchor=south},
  xtick={-3,-1,1,3},xticklabels={,,$1$,$3$},ytick={-2,2},tick label style={font=\footnotesize},clip=false]
\addplot[gray,dashed,domain=-3:3] {x};
\addplot[omThm,thick,samples=2,domain=-3:3] {x/2};
\addplot[omDef,very thick,samples=61,domain=-3:3] {x>1 ? x-1 : (x<-1 ? x+1 : 0)};
\node[font=\footnotesize,anchor=south] at (axis cs:-1,0) {$-1$};
\node[font=\footnotesize,anchor=south] at (axis cs:-3,0) {$-3$};
\node[font=\footnotesize,anchor=west,omDef] at (axis cs:3.05,2) {lasso};
\node[font=\footnotesize,anchor=west,omThm] at (axis cs:3.05,1.5) {ridge};
\node[font=\footnotesize,anchor=west,gray] at (axis cs:3.05,3) {OLS};
\end{axis}
\end{tikzpicture}

\omcaption{The estimate as a function of the OLS coefficient, with $\lambda = 1$ and orthonormal columns. Ridge divides by 2 everywhere. Lasso sets to zero everything inside $[-1, 1]$ and subtracts 1 from the rest, so large coefficients keep a fixed bias.}\label{fig:iv1:shrink}
\end{omfigure}

\iqfollow{Why is there no closed form in general?}{The absolute value is not differentiable, and non-orthogonal columns couple the coefficients. It is solved by coordinate descent, which applies this same rule one coefficient at a time. It can also be written as a quadratic program, by setting $b = b^+ - b^-$ with $b^+, b^- \ge 0$ and using $\sum(b^+ + b^-)$ as the penalty.}
\iqfollow{And without the $\tfrac12$ in front of the squared error?}{The threshold becomes $\lambda/2$.}
\end{solution}

\enlargethispage{-3\baselineskip}%
\begin{solution}{iq:iv1:regression:15}
The coefficients are unbiased under linearity, exogeneity and full rank. The usual standard errors also need errors with constant variance and no correlation. Normal errors are not needed for large samples.

For unbiasedness, you need $y = X\beta + \varepsilon$, $E[\varepsilon \mid X] = 0$ and full column rank. For correct standard errors and the best-linear-unbiased property\index{Gauss--Markov theorem}, you also need homoskedastic and uncorrelated errors. Normality matters only for exact small-sample tests.

On daily stock returns:
\begin{itemize}[leftmargin=*,itemsep=1pt,topsep=2pt]
\item Volatility moves. Coefficients stay unbiased but standard errors are wrong. Use robust (White) standard errors, or scale by a volatility forecast.
\item Errors are correlated, especially with overlapping horizons. Use Newey--West or Hansen--Hodrick standard errors, or non-overlapping data.
\item Fat tails. A few days dominate the fit. The test level holds up reasonably with thousands of points, but check influence, and consider Huber regression or winsorising.
\item Exogeneity fails through omitted factors and look-ahead in features. This one causes bias, and more data will not fix it. Control for factors and time-stamp strictly.
\item Relations drift. Fit on rolling or exponentially weighted windows, and test stability across sub-periods.
\item Regressors are noisy. Slopes shrink towards zero.
\item Regressors are nearly collinear. Coefficients are unstable. Use ridge or combine them.
\end{itemize}

\iqwrong{``The errors must be normal for OLS to be unbiased.'' Normality has nothing to do with bias.}
\end{solution}

\begin{solution}{iq:iv1:regression:16}
The expected $R^2$ is $20/99 \approx 0.20$, the expected adjusted $R^2$ is 0, and about one coefficient has $|t| > 2$.

Under the null, $R^2$ follows a Beta$(p/2,\ (n - p - 1)/2)$ distribution with mean $p/(n - 1)$\index{adjusted $R^2$}. Each of the 20 regressors uses up one of the 99 degrees of freedom beyond the mean, and each one captures an equal share of the variance by chance.

The adjustment $1 - (1 - R^2)(n - 1)/(n - p - 1)$ was built to average zero in this case.

Each $t$-statistic has 79 degrees of freedom, and $P(|t| > 2) \approx 0.049$, so the expected count is $20 \cdot 0.049 \approx 1$.

\iqwhy{The in-sample $R^2$ of 20\% on noise tells you what to expect from a fit with many signals on little data. Compare your results with this baseline, and judge on out-of-sample results.}
\end{solution}

\begin{solution}{iq:iv1:regression:17}
With 250 steps, about 85\% of the regressions are ``significant'', with an average $R^2$ of about $0.24$. With 1,000 steps, the share rises to about 93\%, while $R^2$ stays where it is. Use returns, or test for cointegration.

This is spurious regression\index{spurious regression}. Both series wander, and any two wandering series look related over a finite window. The standard theory assumes stationary errors, and here the residuals are themselves a random walk, so the $t$-statistic has no limiting distribution. It grows like $\sqrt{T}$: going from 250 to 1,000 steps doubles its dispersion. The slope does not converge, and it remains random even with infinite data. These figures come from a seeded simulation.

Differencing cures it. Regressing returns on returns gives $|t| > 1.96$ about 5\% of the time, as it should, with $R^2$ close to zero.

\iqfollow{When is a level regression meaningful?}{When the two series are cointegrated\index{cointegration}, so that a linear combination of them is stationary. The residuals then do not wander, and that is what you test: for example a unit-root test on the residuals, with the appropriate critical values.}
\iqwhy{Two price charts that move together are not evidence. A trading pair needs a stationary spread, and tests on price levels will look wonderful for almost any pair.}
\end{solution}

\begin{solution}{iq:iv1:regression:18}
Ridge, unless you have a reason to believe that a few signals carry almost all of the information.

Lasso assumes sparsity, and it pays for selecting variables with instability: among correlated signals it picks one somewhat at random and changes its pick between samples, which adds turnover. Ridge\index{ridge regression} keeps them all and shrinks along the low-variance directions. With $p/n = 0.4$, many small effects and a weak signal, reducing variance matters most.

In a seeded simulation with 500 signals, correlated through ten factors, and 1,250 observations, where all signals have a small random effect, the out-of-sample $R^2$ is about 1.1\% for ridge and about 0.8\% for lasso, each with its best penalty. With just 5 strong signals and 495 useless ones, lasso comes out ahead.

Change your mind if you have good reason for sparsity, if you want interpretability or fewer inputs to maintain, or if each signal costs money to run. Elastic net is a reasonable compromise.

\iqfollow{What should you be careful of?}{Standardise the signals, do not penalise the intercept, and choose the penalty by time-ordered validation with a gap. Random folds leak information through autocorrelation.}
\end{solution}

\begin{solution}{iq:iv1:regression:19}
The slope is $0.6$, half of $1.2$. Regressing on the current and several lagged market returns and summing the coefficients recovers $1.2$, and a longer return horizon does much the same.

A day's market move shows up in the stock's price on the first later day when it trades, and each day it trades with probability $1/2$. So a fraction $1/2$ of the move shows up the same day, a fraction $1/4$ one day later, $1/8$ two days later, and so on. Regressing on the market's returns at several lags therefore gives slopes $0.6$, $0.3$, $0.15$, and so on, and their sum is the true $1.2$. Ignoring the lags leaves $0.6$.\index{stale prices}

The other cure is to lengthen the return horizon, since a smaller share of the move misses the window. With weekly returns in this model, only the moves late in the week can miss it, and on average $(31/32)/5$ of a week's move does. So the slope is about $1.2\,(1 - 31/160) \approx 0.97$.

\iqwhy{Illiquid names and instruments quoted at different closing times all show low beta and low correlation on daily data. Asian and US closes, for example, are hours apart. Do not mistake this for diversification: at longer horizons the correlation is back. The same effect drives correlation down at very high frequencies.}
\end{solution}

\begin{solution}{iq:iv1:regression:20}
The true standard error is $\sqrt{3/n}$, and the usual formula reports $\sqrt{1/n}$. So it understates the truth by a factor of $\sqrt3 \approx 1.73$.

The slope error is $\sum x_i e_i / \sum x_i^2$, with $e_i = |x_i|\varepsilon_i$. With many observations, its variance is
\[
\frac{E[x^2 e^2]}{n\,(E[x^2])^2} = \frac{E[x^4]}{n} = \frac{3}{n} ,
\]
since $E[x^2 e^2] = E[x^2 \cdot x^2]E[\varepsilon^2]$ and the fourth moment of a standard normal is 3. The usual formula is $s^2/\sum x_i^2$ with $s^2 \to E[e^2] = E[x^2] = 1$, so it gives $1/n$. It uses the average noise, and misses that the noise is largest exactly where $x$ has the most influence on the slope.

To fix the standard errors, use the White sandwich estimator\index{White standard errors}: $\sqrt{\sum x_i^2 \hat e_i^2}/\sum x_i^2$ with demeaned $x$. It estimates the $3/n$ directly and matches the true dispersion in a simulation.

To improve the estimate, use weighted least squares\index{weighted least squares}, which weights each point by the inverse of its noise variance, here $1/x_i^2$. In this case that makes the estimate the mean of $y_i/x_i$, with variance $1/n$, a third of the OLS variance.

\iqfollow{Each of your observations is the average of a different number $m_i$ of underlying data points. What weights do you use?}{$m_i$, since the variance of an average is $\sigma^2/m_i$.}
\iqfollow{You have two independent estimates of the same number, with standard errors 2 and 1. How do you combine them?}{With inverse-variance weights $0.2$ and $0.8$, and a standard error of $\sqrt{0.8} \approx 0.89$.}
\end{solution}

\begin{solution}{iq:iv1:regression:21}
The answers are (a) a factor of 1, so no understatement at all, (b) about $\sqrt{20} \approx 4.47$, and (c) about $3.33$.

The usual formula assumes that $x_t e_t$ is uncorrelated over time. In fact, the variance of the slope is larger by the factor
\[
1 + 2 \sum_{k \ge 1} \rho_x(k)\,\rho_e(k) ,
\]
where $\rho_x$ and $\rho_e$ are the autocorrelations of the regressor and of the error\index{overlapping observations}. Two 20-day windows starting $k$ days apart share $20 - k$ days, so $\rho_e(k) = 1 - k/20$ for $k < 20$.

(a) If the signal is independent from day to day, $\rho_x(k) = 0$, the factor is 1. Overlap in $y$ alone does no harm to the slope's standard error. Many candidates miss this.

(b) If the signal is nearly constant, $\rho_x(k) \approx 1$, the factor is $\sum_{|k|<20}(1 - |k|/20) = 20$, so the standard errors are too small by $\sqrt{20}$.

(c) With $\rho_x(k) = 0.9^k$, the factor is $1 + 2\sum_{k=1}^{19} 0.9^k(1 - k/20) \approx 11.1$, whose square root is $3.33$. Simulations reproduce (a) and (c), and (b) is the limit of the same formula.

The fix is Newey--West\index{Newey--West standard errors} or Hansen--Hodrick standard errors, with the number of lags at least the overlap. Newey--West's weights shrink with the lag, so use more than that: in the simulation with 40 lags it comes within 20\% of the truth for case (c).

\iqfollow{What about plain autocorrelated errors with no overlap?}{With AR(1) in both regressor and error, each with coefficient $0.9$, the factor is $(1 + 0.81)/(1 - 0.81) = 9.53$, so the standard errors are too small by $3.09$. The Durbin--Watson statistic, which is about $2(1 - \rho)$, is the quick check on the residuals.}
\end{solution}

\begin{solution}{iq:iv1:regression:22}
$R^2$ can take any value from $0.36$ to 1. And no, $x_1$ is far from useless: with both regressors the $R^2$ is $0.842$, against $0.16$ for $x_2$ alone.

For two regressors with correlations $r_1$, $r_2$ with $y$ and $\rho$ with each other,
\[
R^2 = \frac{r_1^2 + r_2^2 - 2 r_1 r_2 \rho}{1 - \rho^2} .
\]
With $r_1 = r_2 = 0.6$ this is $0.72/(1 + \rho)$. As $\rho \to 1$, the two regressors are the same variable and $R^2 \to 0.36$. Lowering $\rho$ raises $R^2$. The correlation matrix of $y, x_1, x_2$ must be positive semidefinite, which here requires $\rho \ge 2 \cdot 0.36 - 1 = -0.28$. At that value $R^2 = 1$.

In the second case, with $r_1 = 0$, $r_2 = 0.4$ and $\rho = 0.9$, $R^2 = 0.16/0.19 \approx 0.842$. This is a suppressor variable\index{suppressor variable}: $x_1$ has no relation to $y$ directly, but it cancels the part of $x_2$ that is noise for predicting $y$.

\iqwrong{``Joint $R^2$ is at most the sum of the separate $R^2$ values,'' or ``a variable with zero correlation with $y$ cannot help.'' Both are false.}
\iqwhy{Do not screen signals one at a time by correlation with the target. A hedge is the same idea: an instrument uncorrelated with your target can still reduce your risk once you hold another position.}
\end{solution}

\begin{solution}{iq:iv1:regression:23}
The best combination reaches $0.0577$, $0.0707$ and $0.1$ in the three cases. With correlations $0.05$ and $0.02$ and $\rho = 0.5$, the second signal's weight is negative, even though it predicts positively by itself.

The best linear combination's correlation with the target is $\sqrt{R^2}$\index{signal combination}, and the formula for two regressors gives $R^2 = 2r^2/(1 + \rho)$ with equal correlations. For $r = 0.05$:
\[
\rho = 0.5 : 0.0577, \qquad \rho = 0 : 0.0707, \qquad \rho = -0.5 : 0.1.
\]
Diversification works: correlated signals add little, independent ones add as $\sqrt2$, and negatively correlated signals with the same sign of prediction add more.

The optimal weights are $w \propto \Sigma^{-1}r$. With $\Sigma = \bigl(\begin{smallmatrix}1 & 0.5\\ 0.5 & 1\end{smallmatrix}\bigr)$ and $r = (0.05, 0.02)$,
\[
w = \frac{1}{0.75}\bigl(0.05 - 0.5 \cdot 0.02,\ \ 0.02 - 0.5 \cdot 0.05\bigr) = (0.0533,\ -0.0067).
\]
The second signal's own correlation of $0.02$ is less than the $0.025$ it would get just by being correlated with the first. So it is shorted against the first, which adds a little: the correlation is $0.0503$ against $0.05$ alone.

\iqwhy{This is the regression of return on signals, and it shows why combining beats selecting. The gain from $0.05$ to $0.0577$ is modest but real: an annual Sharpe ratio of about $0.92$, against $0.79$. Note how little it takes for the sign of a weight to flip, so shrink the weights.}
\end{solution}

\begin{solution}{iq:iv1:regression:24}
The best value is $\lambda = \sigma^2/\beta^2 = 4$, which gives a mean squared error of $0.2$ against 1 for OLS, an 80\% reduction.

The ridge estimate $\hat\beta/(1 + \lambda)$ has bias $-\lambda\beta/(1 + \lambda)$ and variance $\sigma^2/(1 + \lambda)^2$, so
\[
\mathrm{MSE}(\lambda) = \frac{\sigma^2 + \lambda^2\beta^2}{(1 + \lambda)^2}.
\]
Setting the derivative to zero gives $\lambda^* = \sigma^2/\beta^2$, and the minimum is $\sigma^2\beta^2/(\sigma^2 + \beta^2)$. With $\sigma = 1$ and $\beta = 0.5$, that is $\lambda^* = 4$ and $\mathrm{MSE} = 0.2$. The estimate is $\hat\beta/5$\index{ridge regression}: throw away four fifths of it.

The weaker the signal relative to noise, the more you should shrink. A signal with a true $t$-statistic of $0.5$ should be cut hard, while one with $\beta = 5$ gets only $\lambda^* = 0.04$. Of course $\beta$ is unknown, which is why you choose $\lambda$ by validation. But the conclusion stands: in low signal-to-noise settings, unbiased estimates are poor bets.

\iqfollow{What is the Bayesian reading?}{With prior $\beta \sim N(0, \tau^2)$, the posterior mean is the ridge estimate with $\lambda = \sigma^2/\tau^2$. For $\tau^2 = 0.25$ that again multiplies the estimate by $1/5$.}
\iqfollow{And with many regressors?}{In the singular value basis, each direction is multiplied by $d^2/(d^2 + \lambda)$, so the directions with little variance in $X$ are shrunk the most.}
\end{solution}

\begin{solution}{iq:iv1:regression:25}
Keeping $x > 0$ leaves the slope at $0.5$ and lowers the correlation to about $0.33$. Keeping $y > 0$ lowers the slope to about $0.22$.

Selecting on $x$ leaves the relation intact\index{selection bias}: given $x$, $y$ still has mean $0.5\,x$, so the slope remains $0.5$. Only the spread of $x$ falls. The variance of a standard normal above zero is $v = 1 - 2/\pi \approx 0.363$. The correlation falls because the explained variance falls:
\[
\text{corr}^2 = \frac{\rho^2 v}{\rho^2 v + 1 - \rho^2} = \frac{0.25 \cdot 0.363}{0.25 \cdot 0.363 + 0.75} \approx 0.108, \quad \text{corr} \approx 0.33.
\]

Selecting on $y$ breaks the relation, because now the noise matters: for low $x$, you keep only the points with a large positive error. So the errors are no longer independent of $x$, and the slope falls. Using the normal relation, with $x = \rho y + \sqrt{1 - \rho^2}\,z$ and the same variance $v$ for $y$ above zero, you get $\mathrm{Cov}(x,y) = \rho v$ and $\mathrm{Var}(x) = \rho^2 v + 1 - \rho^2$, so
\[
\text{slope} = \frac{\rho v}{\rho^2 v + 1 - \rho^2} = \frac{0.5 \cdot 0.363}{0.841} \approx 0.216.
\]

\iqwhy{Filtering on a variable you use as a regressor is safe for the slope. Filtering on the outcome is not, and that includes keeping only survivors, only filled orders or only profitable trades. For example, a study of ``stocks that went up'' or ``orders that filled'' will be biased, and no amount of data will help.}
\end{solution}

\begin{solution}{iq:iv1:regression:26}
Its leverage is $0.101$, removing it lowers the slope by about $0.022$, and its out-of-sample prediction error is about $2.22$.

Leverage\index{leverage} is the diagonal of the hat matrix, how much the observation's own value moves its fitted value. With an intercept,
\[
h = \frac1n + \frac{(x - \bar x)^2}{\sum (x_i - \bar x)^2} = \frac{1}{1000} + \frac{100}{1000} = 0.101.
\]
The leverages sum to the number of coefficients, so the average here is $2/1000$, and this point has over 50 times that.

Deleting a point has exact formulas\index{leave-one-out formula}, so no refit is needed:
\[
\text{prediction error without it} = \frac{e}{1 - h} = \frac{2}{0.899} \approx 2.22,
\]
\[
\text{slope change} = \frac{(x - \bar x)\,e}{(1 - h)\sum(x_i - \bar x)^2} = \frac{10 \cdot 2}{0.899 \cdot 1000} \approx 0.022.
\]
\iqfollow{Why does this matter computationally?}{The same formula gives the full leave-one-out cross-validation error from a single fit, by dividing each residual by $1 - h_i$.}
\iqwhy{One observation with large $x$ can drive the slope, and its residual is small because the line has moved towards it. So look at leverage and not only at residuals. On return data this means the crash days. Check that the sign survives without them, or limit the regressor itself, by winsorising $x$ or using a bounded-influence fit. A Huber loss alone will not do it, since it caps large residuals, and a high-leverage point can have a small one.}
\end{solution}

\begin{solution}{iq:iv1:regression:27}
OLS gives $4/3$, not 1. The instrument gives $\mathrm{Cov}(z, y)/\mathrm{Cov}(z, x) = 1$. A weak instrument gives a very wide estimate, biased towards OLS.

OLS has slope $\mathrm{Cov}(x, y)/\mathrm{Var}(x)$. Here $\mathrm{Var}(x) = 3$, and $\mathrm{Cov}(x, y) = \mathrm{Var}(x) + \mathrm{Cov}(x, u) = 3 + 1 = 4$, so $4/3$: $x$ and the error move together, and the regression gives $x$ the credit.

The variable $z$ is an instrument\index{instrumental variables}: it moves $x$, and it has no other route to $y$. Then $\mathrm{Cov}(z, y) = \beta\,\mathrm{Cov}(z, x)$, since $z$ is uncorrelated with $u$. So $\beta = \mathrm{Cov}(z, y)/\mathrm{Cov}(z, x) = 1/1 = 1$. Two-stage least squares gives the same number: regress $x$ on $z$, then $y$ on the fitted values. But if you do that by hand, the second-stage standard errors are wrong, because they use the wrong residuals.

With a weak instrument, the denominator is close to zero and noisy, and a small error in it moves the ratio a lot. In a simulation with $x = 0.1\,z + u + v$ and 200 observations, the median estimate is pulled well above 1 towards the OLS value, and the spread is very wide. Always look at the first-stage regression.

\iqwhy{A regression of returns on order flow has this problem: prices and flow respond to the same news. Estimating price impact by plain OLS overstates it, since informed flow comes with moves it did not cause. Use flow that is unrelated to news, such as index-rebalance trades or a randomised execution schedule, as an instrument.}
\end{solution}

\begin{solution}{iq:iv1:regression:28}
Expect about $-1\%$, a negative number. You need about 1,980 observations before the out-of-sample $R^2$ reaches zero.

With $p$ estimated coefficients and $n$ observations, each coefficient has an estimation error, and the error of the prediction grows accordingly\index{out-of-sample $R^2$}. With well-behaved regressors, the out-of-sample mean squared error is approximately the noise variance times $(1 + p/n)$. With a total variance of 1, the noise variance is $1 - R^2$, so
\[
R^2_{\text{oos}} \approx R^2 - \frac{p}{n}(1 - R^2) = 0.01 - 0.02 \cdot 0.99 \approx -0.0098 .
\]
The in-sample $R^2$ of the same fit is about 3\%, since it adds $p/n$ rather than subtracting it. Setting the out-of-sample figure to zero gives $n \approx p(1 - R^2)/R^2 = 1980$. For one regressor the break-even is just 99 observations.

\iqwrong{``The true relation is positive, so the fit must beat the mean.'' A real but weak signal, estimated with too many coefficients, loses to the plain average. Its estimation error is larger than its signal.}
\iqwhy{This is the most common reason for a backtest that fails in live trading even with no bug in sight. The fix is to reduce $p$ or to shrink: ridge costs a little bias and removes a large part of the $p/n$ term.}
\end{solution}

\enlargethispage{-2\baselineskip}%
\begin{solution}{iq:iv1:regression:29}
The mean slope is about $+0.03$, even though the truth is 0. The test rejects far too often: about 16\% of the time, not 5\%.

This is the Stambaugh bias\index{Stambaugh bias}. The regressor $x_t$ is uncorrelated with the next shock, but it is not exogenous: it depends on past values of $u$ through $v$. Write $u_t = \gamma v_t + w_t$, with $w$ independent of $v$ and $\gamma = \mathrm{Cov}(u, v)/\mathrm{Var}(v) = -0.9$. The slope of $v_{t+1}$ on $x_t$ is exactly $\hat\rho - \rho$, so
\[
\hat\beta - \beta = \gamma(\hat\rho - \rho) + (\text{slope on } w),
\]
and the last term has mean zero.

So the error in $\hat\beta$ is the error in the persistence estimate, multiplied by $\gamma$. Persistence is underestimated in short samples\index{Kendall bias}: the standard approximation is $E[\hat\rho - \rho] \approx -(1 + 3\rho)/T = -3.85/120 \approx -0.032$. So $E[\hat\beta] \approx (-0.9)(-0.032) \approx 0.029$. A seeded simulation gives a mean slope of $0.032$, a bias in $\hat\rho$ of $-0.035$, and a mean $t$-statistic of about $0.7$. The test rejects about 16\% of the time. You can size this at the board. With $\mathrm{Var}(x) = 1/(1 - \rho^2) \approx 10.3$, the standard error is at least about $1/(3.2\sqrt{120}) \approx 0.029$, as large as the bias, so the $t$-statistic shifts up by at most about 1. On a normal approximation, a shift of $0.7$ to 1 moves the rejection rate from 5\% to between $17\%$ and $26\%$. The simulated 16\% is just below that range, since the mean shift is at the low end and the statistic's spread is a little under 1.

\iqwhy{This is the textbook setting for predicting returns with a valuation ratio such as the dividend yield: a price rise raises returns and lowers the ratio, so $\gamma$ is strongly negative, and the ratio is very persistent. Part of the evidence in such a regression is this bias. Reduce it by adjusting with an estimate of $\gamma(\hat\rho - \rho)$, or test with a null distribution from a simulation.}
\iqfollow{The same bias appears when you fit mean reversion. What does it do?}{It overstates the speed of reversion. With $\rho = 0.95$ the true half-life is $13.5$ periods. The mean estimate on 120 points, $\hat\rho \approx 0.915$, implies $7.8$.}
\end{solution}

\begin{solution}{iq:iv1:regression:30}
The coefficients are about $0.42$ on $w$ and about $0.47$ on $x_2$. You would conclude that $x_2$ has a strong effect, and that is false. Its coefficient only reflects the noise in $w$.

Solve the normal equations in the population\index{errors in variables}. Let $\rho = 0.8$ and $s^2 = 0.5$. Then $\mathrm{Var}(w) = 1 + s^2$, $\mathrm{Cov}(w, x_2) = \rho$, $\mathrm{Cov}(w, y) = 1$ and $\mathrm{Cov}(x_2, y) = \rho$, which gives
\[
\begin{pmatrix} 1 + s^2 & \rho \\ \rho & 1 \end{pmatrix}\begin{pmatrix} b_1 \\ b_2 \end{pmatrix} = \begin{pmatrix} 1 \\ \rho \end{pmatrix}.
\]
The solution is
\[
b_1 = \frac{1 - \rho^2}{1 + s^2 - \rho^2} = \frac{0.36}{0.86} \approx 0.42, \qquad b_2 = \frac{\rho\,s^2}{1 + s^2 - \rho^2} = \frac{0.4}{0.86} \approx 0.47.
\]
When $s^2 = 0$ the coefficient on $x_2$ is zero, as it should be. As $s^2$ grows, the regression trusts $w$ less and leans on $x_2$, because $x_2$ is a clean proxy for $x_1$. In the limit, $b_2$ goes to $\rho = 0.8$, the coefficient of $y$ on $x_2$ alone.

More data does not help. The bias stays while the standard error shrinks, so the $t$-statistic grows without limit. With 2,500 observations it is well above 8 in a simulation.

\iqwrong{Thinking that noise in one regressor affects only that regressor's coefficient. The bias spreads to any correlated regressor.}
\iqwhy{This is how a new factor can ``survive controls'' while adding nothing. If the existing factor you control for is measured with noise, such as a stale book value or a noisy beta, any variable correlated with it will look significant.}
\end{solution}
